%! AP Statistics — Unit 6, Lesson 6.1 — Student Packet Cover Sheet
\documentclass[10pt]{article}
\usepackage{apstats-article}
\usepackage{apstats-boxes}
\usepackage{ltablex}
\keepXColumns

\begin{document}

% Full-bleed navy banner
\begin{tikzpicture}[remember picture, overlay]
  \fill[navy] (current page.north west)
    rectangle ([yshift=-0.9in]current page.north east);
\end{tikzpicture}
\vspace{-0.8in}

\begin{center}
  {\color{white}\textbf{\LARGE Statistics}}\\[4pt]
  {\color{skymid}\large\textbf{Unit 6: Inference for Categorical Data: Proportions}}\\[6pt]
  {\color{white}\textbf{\large Lesson 6.1 \quad Introducing Statistics: Why Be Normal?}}
\end{center}

\vspace{0.22in}
\namedateperiod
\vspace{0.18in}

\begin{learningtargetbox}
\small
\begin{enumerate}[leftmargin=1.5em, itemsep=5pt]
  \item I can identify questions suggested by variation in shapes of distributions of samples
        taken from the same population. \textit{(VAR-1.H)}
  \item I can explain the difference between a confidence interval and a significance test and
        describe when each is appropriate.
\end{enumerate}
\end{learningtargetbox}

\vspace{0.14in}

\begin{tocbox}
% A CONTENTS LIST ONLY: no Score column and no Total row. Whether a component counts is said in
% its description.
\small
\renewcommand{\arraystretch}{1.5}
\begin{tabularx}{\linewidth}{@{} c l X @{}}
\rowcolor{sky}
\textbf{\#} & \textbf{Component} & \textbf{Description} \\
\midrule
1 & Warm-Up        & Sampling distributions review; Large Counts; $SE_{\hat p}$ \\
2 & Group Activity & Inference framing: CI vs.\ significance test scenarios \\
3 & Exit Ticket    & Identify inference goal; verify conditions; compute $SE_{\hat p}$ \\
4 & Homework       & Estimation vs.\ testing classification practice \\
\end{tabularx}
\end{tocbox}

\end{document}
